% =============================================================================
% CAPÍTULO: CONCLUSIÓN
% =============================================================================

El presente trabajo de tesis abordó el diseño, la implementación y la evaluación experimental
de un sistema modular y descentralizado para la optimización del tráfico urbano y la
persistencia verificable de datos asociados. La arquitectura propuesta integra simulación de
tráfico mediante SUMO, clasificación de congestión con inferencia difusa tipo Mamdani,
ajuste adaptativo de temporización semafórica mediante optimización por enjambre de partículas
(PSO) y almacenamiento distribuido basado en IPFS y contratos inteligentes sobre una red
BlockDAG. A diferencia de plataformas centralizadas convencionales, el sistema procesa datos
provenientes de fuentes externas heterogéneas, expone sus capacidades mediante APIs REST y
garantiza trazabilidad e integridad de los registros sin depender de un operador único de
confianza.

Los capítulos anteriores desarrollaron los fundamentos teóricos, el diseño de cada módulo y
su interacción dentro del flujo completo: desde la generación de escenarios controlados y la
detección de cuellos de botella en traffic-sim, pasando por el análisis y la
optimización en traffic-sync, hasta la publicación auditable de resultados en
traffic-storage. El capítulo de evaluación experimental consolidó estas piezas con
mediciones reproducibles en hardware de consumo,
demostrando que el sistema es viable para entornos municipales con recursos limitados.

\section{Resultados principales}

La evaluación del módulo traffic-sync confirmó una escalabilidad aproximadamente
lineal del tiempo de procesamiento conforme aumenta el número de sensores simulados: para
2000 sensores, el ciclo completo de análisis y optimización se completó en 67,6~s
($\approx 0{,}034$~s por sensor). El consumo de memoria RAM se mantuvo estable entre
198 y 219~MB en todos los escenarios, y el uso de CPU disminuyó del 72,3\% al 50,3\% a
medida que creció la concurrencia, comportamiento coherente con el procesamiento asíncrono
de FastAPI y con el bajo costo computacional unitario de los componentes difuso y PSO.
Estos resultados respaldan la viabilidad de despliegues sobre hardware embebido de bajo costo,
como unidades Raspberry~Pi en campo.

En traffic-storage, las transacciones de almacenamiento alcanzaron confirmación en
la red BlockDAG en $\approx$1~s, con costos inferiores a \$1{,}2$\times$10\textsuperscript{-5}~USD
por transacción. Los tiempos de subida extremo a extremo se situaron entre 6,7 y 7,3~s según
la densidad vehicular del escenario, mientras que la descarga se completó en $\approx$1,75~s,
adecuado para flujos de auditoría y consulta sobre datos agregados de ciclos completados, no
sobre control en tiempo estricto. La verificación de integridad mediante el CID de IPFS
demostró que los contenidos recuperados conservan fielmente el original. Frente a Ethereum,
el sistema superó en más del 95\% la latencia de confirmación y en más de cuatro órdenes de
magnitud el costo operativo; frente a arquitecturas IoT-cloud centralizadas, la reducción de
latencia de persistencia verificable superó el 60\%, con inmutabilidad garantizada por
consenso distribuido y no por confianza en un operador central.

En conjunto, los experimentos establecen que la combinación IPFS+BlockDAG constituye una
alternativa técnicamente competitiva para monitoreo urbano de alta frecuencia con requisitos
de trazabilidad auditables, mientras que el núcleo analítico escala de forma predecible hacia
escenarios de optimización por zonas urbanas.

\section{Divulgación científica}

Los hallazgos centrales de esta investigación fueron sintetizados y publicados en el
artículo de conferencia internacional titulado \emph{Scalable and Decentralized Urban Traffic
Optimization with Verifiable Storage for Smart City Deployments},
presentado en la \textit{2025 IEEE CHILEAN Conference on Electrical, Electronics Engineering,
Information and Communication Technologies} (CHILECON), celebrada en Valparaíso, Chile, en
octubre de 2025.\footnote{Disponible en IEEE Xplore:
\nolinkurl{https://ieeexplore.ieee.org/document/11476672}.}
El artículo documenta, en formato condensado,
la arquitectura modular del sistema, las métricas de escalabilidad y almacenamiento aquí
ampliadas, y la comparación con plataformas basadas en Ethereum y soluciones centralizadas.
La publicación valida externamente los resultados experimentales de la tesis y posiciona la
propuesta en el ámbito de sistemas inteligentes para ciudades con almacenamiento verificable
y despliegue en dispositivos de bajo costo.

\section{Limitaciones}

Pese a los resultados favorables, deben reconocerse limitaciones relevantes para un
despliegue productivo. En primer lugar, las evaluaciones se realizaron principalmente en
entorno simulado y sobre hardware de laboratorio; la transferencia a escenarios urbanos reales
implica heterogeneidad de sensores, variabilidad en las comunicaciones y resistencia
institucional a soluciones descentralizadas, aspectos aún no validados en pilotos de campo.

En segundo lugar, el módulo traffic-sync procesa actualmente un sensor por solicitud;
si bien la arquitectura y las pruebas de estrés anticipan la evaluación conjunta por zonas,
la optimización multivariable por clústeres de semáforos permanece como trabajo pendiente para
lograr una dispersión efectiva del tráfico a escala territorial.

En tercer lugar, traffic-storage dependió de un nodo IPFS gestionado por terceros
(Pinata), introduce variabilidad en los tiempos de minado de la red pública de pruebas
BlockDAG y concentra el cuello de botella en la confirmación on-chain ($\approx$3,5~s dentro
del pipeline de subida). No se implementaron aún mecanismos activos de tolerancia a fallos ni
replicación automática de contenido.

Finalmente, la configuración operativa de las APIs requiere validación exhaustiva del entorno
antes del despliegue, como evidenció el incidente de redirecciones HTTP~307 durante las primeras
corridas del protocolo de prueba.

\section{Trabajo futuro}

Como continuación natural de este trabajo, se identifican las siguientes líneas de investigación
y desarrollo:

\begin{itemize}
  \item Implementar agrupamiento jerárquico de sensores y optimización PSO por zonas urbanas,
        habilitando la coordinación de clústeres de semáforos en paralelo.
  \item Integrar fuentes de datos reales provenientes de dispositivos de conteo, sensores IoT
        con conectividad LoRaWAN y sistemas de visión por computadora.
  \item Desplegar nodos IPFS propios y evaluar redes BlockDAG privadas o permisionadas para
        despliegues municipales, reduciendo la variabilidad de confirmación y la dependencia
        de proveedores externos.
  \item Optimizar la gestión de nonces y la precarga del estado de gas para acortar la fase de
        confirmación en el pipeline de almacenamiento.
  \item Diseñar esquemas de comunicación asíncrona mediante instancias intermedias que alimenten
        en tiempo real los módulos del sistema desde el campo.
  \item Ejecutar pilotos en ciudades intermedias de la región para medir resiliencia, calidad de
        servicio e impacto en la fluidez del tráfico frente a plataformas centralizadas.
  \item Explorar la integración con marcos de \textit{digital twins} urbanos para planificación
        y simulación avanzada.
\end{itemize}

\section{Síntesis final}

El sistema desarrollado en esta tesis ofrece una base tecnológica sólida que responde a
criterios de trazabilidad, seguridad, integridad y eficiencia operativa en la gestión del
tráfico urbano. Su enfoque modular y descentralizado, respaldado por evidencia cuantitativa y
por la publicación derivada en CHILECON~2025, lo posiciona como una propuesta
viable para municipios con presupuesto limitado. Las mejoras futuras deberán cerrar la brecha
entre los resultados de laboratorio y la evidencia práctica en entornos urbanos reales,
fortaleciendo la tolerancia a fallos y la optimización territorial del control semafórico.
